\chapter{The ability population: hierarchical empirical Bayes}
\label{ch:pop}

\section{The shrinkage problem}

Chapter~\ref{ch:energetics} produced, for each run of each horse, a noisy pair
$(\widehat\vc,\widehat\Wp)$. We need a per-horse \emph{estimate} of the latent ability to feed the
drift, and it must satisfy three requirements simultaneously: it must (i) borrow strength across
horses so a lightly-raced horse is not at the mercy of one fluke run, (ii) condition on the
horse's class and trip so the prior it shrinks toward is appropriate, and (iii) be strictly
leak-free for walk-forward backtesting. The classical tool is a one-way Normal--Normal
random-effects model fitted by empirical Bayes, with a rating-conditioned fixed effect; this
chapter develops it for each of $\vc$ and $\Wp$ independently.

\section{The Normal--Normal random-effects model}

Work with one parameter at a time (say $\vc$; $\Wp$ is identical). Let $y_{hr}$ denote the per-run
estimate \eqref{eq:vc-est} for the $r$-th run of horse $h$. Model it as a fixed effect plus a
horse random effect plus within-horse noise:
\begin{equation}
  y_{hr} \;=\; \underbrace{f(b_{hr}, R_{hr})}_{\text{fixed effect}}
            \;+\; \underbrace{\theta_h}_{\text{horse effect}}
            \;+\; \underbrace{\varepsilon_{hr}}_{\text{within-horse noise}},
  \qquad
  \theta_h \sim \Norm(0,\tau^2), \quad \varepsilon_{hr}\sim\Norm(0,\sigma^2).
  \label{eq:randeff}
\end{equation}
Here $b_{hr}\in\{\text{sprint},\text{mile},\text{middle},\text{staying}\}$ is the distance band
of the run and $R_{hr}$ is the horse's official rating going into it. The two variance components
have direct interpretations: $\sigma^2$ is the \emph{within-horse} run-to-run variance (how noisy
a single sectional estimate is), and $\tau^2$ is the \emph{between-horse} variance (how much
genuine spread there is in true ability after conditioning on class and trip).

\subsection{The rating-conditioned fixed effect}

The fixed effect is not a single grand mean but a per-band intercept plus a global rating slope,
\begin{equation}
  f(b,R) \;=\; \beta_b \;+\; \gamma\,(R-\bar R),
  \label{eq:fixed}
\end{equation}
so that a horse's prior centres on the value expected \emph{for its class and trip}, not on the
global average. The band intercepts $\beta_b$ capture that $\vc$ falls with distance (a sprinter's
sustainable velocity exceeds a stayer's); the slope $\gamma$ captures that a higher-rated horse
has, on average, a higher critical velocity. The slope is estimated by an ordinary least-squares
fit of the band-centred residual on the centred rating,
\begin{equation}
  \widehat\gamma
   \;=\; \frac{\sum_{hr}(R_{hr}-\bar R)\,(y_{hr}-\beta_{b_{hr}})}{\sum_{hr}(R_{hr}-\bar R)^2},
  \label{eq:slope}
\end{equation}
after which the rating-explained part is removed to form the pure horse-plus-noise residual
$u_{hr} = y_{hr} - \beta_{b_{hr}} - \widehat\gamma(R_{hr}-\bar R)$. Conditioning on rating is what
makes the population a \emph{segmented-by-class} prior rather than a mono-population; it matters
because a debut or lightly-raced horse shrinks almost entirely to this prior, so the prior had
better know the horse's class.

\subsection{Estimating the variance components (ANOVA method of moments)}

The two variance components are estimated by the ANOVA method of moments --- robust, closed-form,
and assumption-light. The within-horse variance is the pooled sample variance over horses with at
least two runs,
\begin{equation}
  \widehat{\sigma}^2
   \;=\; \frac{\sum_{h:\,n_h\ge 2}(n_h-1)\,s_h^2}{\sum_{h:\,n_h\ge 2}(n_h-1)},
  \qquad
  s_h^2 = \frac{1}{n_h-1}\sum_{r}(u_{hr}-\bar u_{h})^2,
  \label{eq:sigma2}
\end{equation}
where $n_h$ is horse $h$'s run count and $\bar u_h$ its mean residual. The between-horse component
follows from equating the between-group mean square to its expectation. With
$\mathrm{SSB}=\sum_h n_h \bar u_h^2$ (the residuals are already fixed-effect-centred so the grand
mean is zero), $H$ horses, and total runs $n_\bullet=\sum_h n_h$,
\begin{equation}
  \mathrm{MSB} = \frac{\mathrm{SSB}}{H-1},
  \qquad
  n_0 = \frac{1}{H-1}\Bigl(n_\bullet - \frac{\sum_h n_h^2}{n_\bullet}\Bigr),
  \qquad
  \widehat{\tau}^2 = \max\!\Bigl(\frac{\mathrm{MSB}-\widehat\sigma^2}{n_0},\ \epsilon\Bigr),
  \label{eq:tau2}
\end{equation}
where $n_0$ is the usual effective group size for unbalanced designs and the floor $\epsilon$
keeps $\widehat\tau^2$ strictly positive. (The implementation, \code{calibration.py}, computes
exactly these quantities in vectorised \code{polars}/\code{numpy}.)

\section{The partially-pooled per-horse estimate}

With the components in hand, the empirical-Bayes posterior mean for horse $h$ is the
fixed-effect prediction plus a \emph{shrunken} version of its own mean residual,
\begin{equation}
  \boxed{\;\widehat{\vc}_h \;=\; f(b,R) \;+\; \lambda(n_h)\,\bar u_h,\qquad
  \lambda(n_h) \;=\; \frac{\tau^2}{\tau^2 + \sigma^2/n_h}\;}
  \label{eq:shrink}
\end{equation}
The shrinkage weight $\lambda(n_h)\in[0,1)$ is the heart of the model. Read its limits:
\begin{itemize}
  \item $n_h\to 0$ (debut): $\lambda\to 0$, the estimate is the pure rating-conditioned population
        mean --- the horse has provided no evidence, so it is its class.
  \item $n_h\to\infty$ (established): $\lambda\to 1$, the estimate trusts the horse's own runs
        almost entirely.
  \item Intermediate $n_h$: $\lambda$ interpolates, and the cross-over is governed by the
        \emph{noise-to-signal ratio} $\sigma^2/\tau^2$: a noisier measurement (large $\sigma^2$)
        or a more homogeneous population (small $\tau^2$) means more pooling for the same run
        count.
\end{itemize}
The posterior also has a variance, used later by the precision-fusion of Chapter~\ref{ch:fusion}:
\begin{equation}
  \mathrm{sd}(\widehat{\vc}_h)
   \;=\; \Bigl(\frac{n_h}{\sigma^2} + \frac{1}{\tau^2}\Bigr)^{-1/2},
  \label{eq:postsd}
\end{equation}
the standard ``precisions add'' result: the posterior precision is the prior precision $1/\tau^2$
plus $n_h$ run-precisions $1/\sigma^2$. This identity --- that each independent source of evidence
contributes its precision additively --- is the conceptual seed of the entire fusion architecture.

\section{Leak-free variants: leave-one-out and strictly-prior}

The shrinkage formula \eqref{eq:shrink} must be made leak-free in two distinct settings, and the
implementation provides both.

\paragraph{Leave-one-out (for in-sample calibration targets).} When the energetics of a run are
used to \emph{calibrate} against that same run's outcome, the run must be excluded from its own
horse's pool. The LOO estimate replaces $\bar u_h$ and $n_h$ by their values with the current run
dropped,
\begin{equation}
  \bar u_{h}^{(-r)} \;=\; \frac{\sum_{r'\ne r} u_{hr'}}{n_h-1},
  \qquad
  \lambda(n_h-1),
  \label{eq:loo}
\end{equation}
computed in vectorised form as (group sum $-$ this run)/(group count $-1$). This is the quantity a
race model may consume without a run's own outcome leaking into the fit that predicts it.

\paragraph{Strictly-prior expanding window (for walk-forward forecasting).} For a genuine
forecast of a race on date $d$, the estimate may use only runs \emph{strictly before} $d$. The
trailing variant sorts each horse's runs by date and forms an expanding mean over earlier runs
only,
\begin{equation}
  \widehat{\vc}_h(d) \;=\; f(b,R) \;+\; \lambda(n_{<d})\,
   \frac{1}{n_{<d}}\!\!\sum_{r:\,\text{date}(r)<d}\!\! u_{hr},
  \label{eq:trailing}
\end{equation}
with $n_{<d}$ the count of prior runs (and a fallback to the pure population mean when
$n_{<d}=0$). The population parameters $(\beta_b,\gamma,\sigma^2,\tau^2)$ are themselves fit on the
training window only, so the whole pipeline --- population and per-horse estimate --- is leak-free.
This trailing estimate is the one the walk-forward backtest of Chapter~\ref{ch:method} actually
consumes.

\section{Output of the population layer}

The population layer delivers, for every runner in every race, a leak-free
$(\widehat\vc_h,\widehat\Wp_h)$ together with the horse's prior-run count $n_h$ (which the
shrinkage and the later fusion both need) and the population's variance components. These are the
\emph{baseline} drift and reserve --- the ``rating-conditioned mean plus shrunken own-history''
that a horse carries into a race before any of the feature channels of the next chapters nudge it.
In the assembled drift of Chapter~\ref{ch:assembly} this appears as the leading terms
$\code{pop.v\_c.fixed}(b,R) + \texttt{SHRINK\_MULT}\cdot\lambda(n)\,u_{\vc}$ and the analogous
reserve term, on top of which every ability channel adds.
